\section{Introduction}
Après avoir présenté le cadre théorique et méthodologique de cette étude, ce chapitre est consacré à la mise en œuvre de l’expérimentation ainsi qu’à son suivi tout au long du cycle cultural. Il décrit les différentes étapes de la conduite de l’essai, depuis l’aménagement des parcelles et l’installation des cultures jusqu’au suivi de leur évolution. Une attention particulière est accordée au dispositif d’irrigation, aux outils et méthodes de collecte des données, ainsi qu’à la gestion hebdomadaire des parcelles. Les approches déployées pour assurer la protection des cultures contre les maladies et les ravageurs sont également présentées. Finalement, l’évolution des divers plots au cours de l’expérimentation est illustrée afin de mettre en évidence les circonstances dans lesquelles les données ont été recueillies. L'ensemble de ces éléments constitue une base essentielle à la compréhension et à l'interprétation des résultats présentés dans le chapitre suivant.